\subsection{\texorpdfstring{函数 \(e^{z}\) 的性质}{函数 e 的 z 次幂的性质}}

\(z \mapsto e^z\) 是 \(\mathbf{C}\) 到 \(\mathbf{C}^*\) 上的同态，这一事实可用下列恒等式表达

\[
e ^ {z + z ^ {\prime}} = e ^ {z} e ^ {z ^ {\prime}}, \qquad e ^ {0} = 1, \qquad e ^ {- z} = 1 / e ^ {z}.\tag{22}
\]

按定义，对每个 \(z = x + iy\) ，有

\[
e ^ {x + i y} = e ^ {x} (\cos y + i \sin y)\tag{23}
\]

并且由于 \(e^x > 0\) ，可见 \(e^z\) 的绝对值为 \(e^x\) ，辐角为 \(y\) （模 \(2\pi\) ）。

特别地，定义 2（III, p.~98）给出

\[
\mathbf {e} (x) = e ^ {2 \pi i x}\tag{24}
\]

它使我们可以把定义 \(\cos x\) 与 \(\sin x\) 的公式写成如下形式

\[
\cos x = \frac {1}{2} \left(e ^ {i x} + e ^ {- i x}\right), \quad \sin x = \frac {1}{2 i} \left(e ^ {i x} - e ^ {- i x}\right)\tag{25}
\]

（欧拉公式）。

由于 \(2\pi\) 是 \(\mathbf{e}(y/2\pi)\) 的最小周期，\(2\pi i\) 是 \(e^{z}\) 的最小周期；换句话说，\(e^{z}\) 的周期群是数 \(2n\pi i\) 所成的集，其中 n 跑遍 Z。

最后，III, p.~98 的公式 (21) 可写成

\[
\mathrm{D} \left(e ^ {z}\right) = e ^ {z}\tag{26}
\]

从而，对每个复数 \(a\)

\[
\mathrm{D} \big (e ^ {a z} \big) = a e ^ {a z}.\tag{27}
\]

注。若在公式 (27) 中把函数 \(e^{az}\) （a 为复数）限制在实轴上，则对实数 x 再次得到

\[
\mathrm{D} (e ^ {a x}) = a e ^ {a x}.\tag{28}
\]

这个公式使我们可以对 \(e^{\alpha x} \cos \beta x\) ，\(e^{\alpha x} \sin \beta x\) （\(\alpha\) 与 \(\beta\) 为实数）中的每一个函数求出原函数；事实上，我们有 \(e^{(\alpha + i\beta)x} = e^{\alpha x} \cos \beta x + ie^{\alpha x} \sin \beta x\) ，于是由 (28)

\[
\begin{array}{l} \mathrm{D} \left(\mathcal {R} \left(\frac {1}{\alpha + i \beta} e ^ {(\alpha + i \beta) x}\right)\right) = e ^ {\alpha x} \cos \beta x \\ \mathrm{D} \left(\mathcal {I} \left(\frac {1}{\alpha + i \beta} e ^ {(\alpha + i \beta) x}\right)\right) = e ^ {\alpha x} \sin \beta x. \end{array}
\]

同样，可以把求 \(x^n e^{\alpha x} \cos \beta x\) 或 \(x^n e^{\alpha x} \sin \beta x\) （ \(n\) 为整数 \(> 0\) ）的原函数归结为求 \(x^n e^{(\alpha + i\beta)x}\) 的原函数；而 \(n + 1\) 阶分部积分公式（II, p.~60，公式 (11)）表明，后一函数的一个原函数是

\[
e ^ {(\alpha + i \beta) x} \left[ \frac {x ^ {n}}{\alpha + i \beta} - \frac {n x ^ {n - 1}}{(\alpha + i \beta) ^ {2}} + \frac {n (n - 1) x ^ {n - 2}}{(\alpha + i \beta) ^ {3}} + \dots + (- 1) ^ {n} \frac {n !}{(\alpha + i \beta) ^ {n + 1}} \right].
\]

另一方面，利用欧拉公式可以把 \(\cos x\) 或 \(\sin x\) 的每个正整数次幂表示为指数函数 \(e^{ipx}\) （ \(p\) 为正或负整数）的线性组合。于是由公式 (28)，可以把形如 \(x^n e^{\alpha x}(\cos \beta x)^r (\sin \gamma x)^s\) 的函数的原函数表示成形如 \(x^p e^{\alpha x}\cos \lambda x\) 与 \(x^p e^{\alpha x}\sin \mu x\) 的函数的线性组合（\((n, p, r, s\) 为整数，\(\alpha, \beta, \gamma, \lambda, \mu\) 为实数）。

例。有

\[
\sin^ {2 n} x = \frac {(- 1) ^ {n}}{2 ^ {2 n}} \left(e ^ {i x} - e ^ {- i x}\right) ^ {2 n} = \frac {(- 1) ^ {n}}{2 ^ {2 n}} \left(e ^ {2 n i x} - \binom {2 n} {1} e ^ {(2 n - 2) i x} + \dots + e ^ {- 2 n i x}\right)
\]

由此得

\[
\begin{array}{l} \int_ {0} ^ {x} \sin^ {2 n} t d t = \frac {(- 1) ^ {n}}{2 ^ {2 n}} \left(\frac {1}{n} \sin 2 n x - \binom {2 n} {1} \frac {1}{n - 1} \sin (2 n - 2) x + \dots \right. \\ \qquad \qquad \qquad \qquad \qquad \qquad \qquad \qquad \qquad \qquad + (- 1) ^ {n - 1} \binom {2 n} {n - 1} \sin 2 x + (- 1) ^ {n} \binom {2 n} {n} x \Bigg) \end{array}
\]

特别地，

\[
\int_ {0} ^ {\pi / 2} \sin^ {2 n} t   d t = \binom {2 n} {n} \frac {1}{2 ^ {2 n}}   \frac {\pi}{2} = \frac {1 . 3 . 5 \ldots (2 n - 1)}{2 . 4 . 6 \ldots 2 n}   \frac {\pi}{2}.\tag{29}
\]

